\section{当下共识与竞争：国内团队如何追赶}

罗福莉认为，Anthropic 的路径正确已经成为共识；在路径更清晰后，国内大模型团队进入加速追赶状态。现在 pre-train 代差很小，甚至国内在某些结构上有优势；真正要 all in 的，是 Agent post-train，尤其是在 Agent 上怎么做好 RL scaling。

\subsection{Coding 为什么每个范式都重要}

Coding 在多个阶段都戳中了关键点。预训练阶段，代码数据质量高、结构清晰；reasoning 阶段，代码和数学有可验证反馈；Agent 阶段，软件开发本身是长程任务，有真实环境、测试、调试和复杂依赖。它既是训练信号，也是产品场景。

\begin{importantbox}{Coding 的三重角色}
Coding 是数据、评估和产品。作为数据，它结构化且质量高；作为评估，它有测试和运行反馈；作为产品，它直接对应高价值生产力场景。这就是为什么 coding 在每个范式中都重要。
\end{importantbox}

\subsection{环境比经验更重要}

结尾部分，罗福莉强调环境比经验更重要。所谓环境，包括算力、基础设施、芯片、操作系统、流量、社交、组织文化和战略资源。一个团队是否能在新范式中领先，取决于它能否让这些资源合力适配 Agent 架构。

\begin{warningbox}{经验可能成为负担}
在范式快速变化时，旧经验会帮助判断，也会限制想象。真正重要的是环境能否给团队持续反馈和加速度，使其能不断否认旧判断、重构旧流程。
\end{warningbox}

\subsection{本章小结}

当下竞争的共识是：pre-train 差距缩小，Agent post-train 成为主线。国内团队的机会在于路径更清晰、成本敏感、应用反馈快；挑战在于 RL infra、任务环境和组织重构。

